\documentclass{exam}
\usepackage[utf8]{inputenc}
\usepackage[T1]{fontenc}
\usepackage{geometry}
\usepackage{graphicx}
\usepackage{float}
\usepackage{hyperref}
\usepackage[english,indonesian]{babel}

\title{Optimasi Pemrosesan Template LaTeX Menggunakan Pendekatan Berbasis AI}
\author{Academic Support Agent}
\date{\today}

\begin{document}
\maketitle

\begin{abstract}

\end{abstract}

\chapter{PENDAHULUAN}


\chapter{TINJAUAN PUSTAKA}


% Original content (commented out):
% \documentclass[12pt, addpoints]{exam}
\usepackage{mayencourt}
\pgfplotsset{compat=1.15}
%\usepgfplotslibrary{statistics}
%\usetikzlibrary{arrows}
\usetikzlibrary{patterns}

\pagestyle{headandfoot}
\extraheadheight{1cm}
\runningheadrule
\firstpageheader{\includegraphics[scale=1]{logo.png}}
                {    \textbf{ Maturité professionnelle, orientation : \\
                      Economie et services, type "Economie"\\
                      Branche :\\
                      Mathématiques}
                 }
                {   \textbf{Formative\\
                    2018-2019\\
                    Série A\\
                    Durée 120 minutes}}
\runningheader{\includegraphics[scale=1]{logo.png}}
                {    \textbf{ Maturité professionnelle, orientation : \\
                      Economie et services, type "Economie"\\
                      Branche :\\
                      Mathématiques}
                 }
                {   \textbf{Formative\\
                    2018-2019\\
                    Série A\\
                    Durée 120 minutes}}




\begin{document}
\printanswers

\vspace{5cm}

\begin{center}
\begin{tabular}{|l|l|}
\hline
&\\
\makebox[0.4\textwidth]{Nom : \enspace\hrulefill}    & \makebox[0.4\textwidth]{N candidat : \enspace\hrulefill}  \\
&\\
\makebox[0.4\textwidth]{Prénom : \enspace\hrulefill}     & 
\makebox[0.4\textwidth]{Date : \enspace\hrulefill}\\
&\\
\hline
\end{tabular}
\end{center}

\vspace{5cm}

\begin{center}
\gradetable[h][questions]
\end{center}

\vspace{2cm}

\begin{flushright}
\begin{tabular}{|l|p{2cm}|}
\hline
     &  \\
    Note & \\
     & \\
\hline
\end{tabular}
\end{flushright}

\newpage

\textbf{Moyen auxiliaire autorisé :} calculatrice non graphique, non programmable.

\vspace{3cm}

\textbf{Remarques :}
\begin{itemize}
    \item Tous les calculs et développements des solutions doivent figurer sur les feuilles de données au propre. Les réponses sans développement mathématique ne sont pas prises en compte. Un espace réponse supplémentaire est disponible en pages 14 et 15.
    \item Les développements doivent être suffisamment détaillés pour que l'on puisse suivre le raisonnement et les opérations aboutissant au résultat.
    \item Les réponses des exercices seront simplifiées au maximum.
\end{itemize}

\vspace{5cm}

\begin{tabular}{ll}
Signature des experts : & \makebox[0.4\textwidth]{ \enspace\hrulefill}	\\
&\\
&\\
& \makebox[0.4\textwidth]{ \enspace\hrulefill}\\
\end{tabular}

\newpage

\begin{questions}

\question Calculez et simplifiez
\begin{parts}
\part $$ \frac{3}{5}\div 12 - \left( \frac{3}{4}-\frac{7}{6}\right) \cdot \frac{9}{2}$$
\begin{tikzpicture}
	% Dimensions du repere
	\def\xmin{0} \def\xmax{15} \def\ymin{0} \def\ymax{5}
	% Grilles
%	\draw [step=0.1cm,gray,ultra thin]  (\xmin,\ymin) grid (\xmax,\ymax);
	\draw [step=0.5cm,gray, ultra thin] (\xmin,\ymin) grid (\xmax,\ymax);
	\draw [step=1cm,gray, thin] (\xmin,\ymin) grid (\xmax,\ymax);
	\draw [step=5cm,gray,thick] (\xmin,\ymin) grid (\xmax,\ymax);
\end{tikzpicture}
\part $$ \sqrt{108} -3\cdot \sqrt{8} \cdot \left(\sqrt{27}-\sqrt{25}\right)+\sqrt{24}$$
\begin{tikzpicture}
	% Dimensions du repere
	\def\xmin{0} \def\xmax{15} \def\ymin{0} \def\ymax{5}
	% Grilles
%	\draw [step=0.1cm,gray,ultra thin]  (\xmin,\ymin) grid (\xmax,\ymax);
	\draw [step=0.5cm,gray, ultra thin] (\xmin,\ymin) grid (\xmax,\ymax);
	\draw [step=1cm,gray, thin] (\xmin,\ymin) grid (\xmax,\ymax);
	\draw [step=5cm,gray,thick] (\xmin,\ymin) grid (\xmax,\ymax);
\end{tikzpicture}
\part $$ \frac{x^7\cdot (2x)^2 \cdot \left(6 x^0\right)^{-3}}{x^2 \cdot (-2x)^{-5}} \div \frac{(3x)^2 \cdot (9x)^3}{(6x)^4}$$
\begin{tikzpicture}
	% Dimensions du repere
	\def\xmin{0} \def\xmax{15} \def\ymin{0} \def\ymax{5}
	% Grilles
%	\draw [step=0.1cm,gray,ultra thin]  (\xmin,\ymin) grid (\xmax,\ymax);
	\draw [step=0.5cm,gray, ultra thin] (\xmin,\ymin) grid (\xmax,\ymax);
	\draw [step=1cm,gray, thin] (\xmin,\ymin) grid (\xmax,\ymax);
	\draw [step=5cm,gray,thick] (\xmin,\ymin) grid (\xmax,\ymax);
\end{tikzpicture}
\end{parts}

\question Résolvez :
\begin{parts}
\part $$ \left\{
\begin{array}{lcl}
    y-3\cdot  \displaystyle\frac{2x-1}{6} &=& \frac{1}{2}\cdot \left(1+ \displaystyle \frac{x}{3} \right) \\
    &\\
    2x-3y + 7 &=&  \displaystyle \frac{x+2}{3}-\frac{y-1}{5}
\end{array} 
\right.$$
\begin{tikzpicture}
	% Dimensions du repere
	\def\xmin{0} \def\xmax{15} \def\ymin{0} \def\ymax{8}
	% Grilles
%	\draw [step=0.1cm,gray,ultra thin]  (\xmin,\ymin) grid (\xmax,\ymax);
	\draw [step=0.5cm,gray, ultra thin] (\xmin,\ymin) grid (\xmax,\ymax);
	\draw [step=1cm,gray, thin] (\xmin,\ymin) grid (\xmax,\ymax);
	\draw [step=5cm,gray,thick] (\xmin,\ymin) grid (\xmax,\ymax);
\end{tikzpicture}
\part $$3x-5 \cdot (x+2) \geq 2x+4-\left(\frac{x+3}{2}\right)$$
\begin{tikzpicture}
	% Dimensions du repere
	\def\xmin{0} \def\xmax{15} \def\ymin{0} \def\ymax{8}
	% Grilles
%	\draw [step=0.1cm,gray,ultra thin]  (\xmin,\ymin) grid (\xmax,\ymax);
	\draw [step=0.5cm,gray, ultra thin] (\xmin,\ymin) grid (\xmax,\ymax);
	\draw [step=1cm,gray, thin] (\xmin,\ymin) grid (\xmax,\ymax);
	\draw [step=5cm,gray,thick] (\xmin,\ymin) grid (\xmax,\ymax);
\end{tikzpicture} 
\end{parts}

\question Soit la parabole d’équation $y=-2 x^2 + 3 x + 2$ et la droite d’équation $y=3x$.

\begin{parts}
\part Calculez tous les points particuliers de la parabole (sommet, intersections avec les axes).

\begin{tikzpicture}
	% Dimensions du repere
	\def\xmin{0} \def\xmax{15} \def\ymin{0} \def\ymax{9}
	% Grilles
%	\draw [step=0.1cm,gray,ultra thin]  (\xmin,\ymin) grid (\xmax,\ymax);
	\draw [step=0.5cm,gray, ultra thin] (\xmin,\ymin) grid (\xmax,\ymax);
	\draw [step=1cm,gray, thin] (\xmin,\ymin) grid (\xmax,\ymax);
	\draw [step=5cm,gray,thick] (\xmin,\ymin) grid (\xmax,\ymax);
\end{tikzpicture}
\part Calculez les points d’intersection de ces deux courbes.

\begin{tikzpicture}
	% Dimensions du repere
	\def\xmin{0} \def\xmax{15} \def\ymin{0} \def\ymax{9}
	% Grilles
%	\draw [step=0.1cm,gray,ultra thin]  (\xmin,\ymin) grid (\xmax,\ymax);
	\draw [step=0.5cm,gray, ultra thin] (\xmin,\ymin) grid (\xmax,\ymax);
	\draw [step=1cm,gray, thin] (\xmin,\ymin) grid (\xmax,\ymax);
	\draw [step=5cm,gray,thick] (\xmin,\ymin) grid (\xmax,\ymax);
\end{tikzpicture} 
\newpage
\part c)	Représentez graphiquement les deux courbes en faisant apparaître tous les points calculés. (unité : 2 carrés = 1)

\begin{tikzpicture}
	% Dimensions du repere
	\def\xmin{0} \def\xmax{15} \def\ymin{0} \def\ymax{20}
	% Grilles
%	\draw [step=0.1cm,gray,ultra thin]  (\xmin,\ymin) grid (\xmax,\ymax);
	\draw [step=0.5cm,gray, ultra thin] (\xmin,\ymin) grid (\xmax,\ymax);
	\draw [step=1cm,gray, thin] (\xmin,\ymin) grid (\xmax,\ymax);
	\draw [step=5cm,gray,thick] (\xmin,\ymin) grid (\xmax,\ymax);
\end{tikzpicture} 
\end{parts}

\newpage

\question Par rapport à la liste des vitesses des voitures relevée devant une école, construire le tableau complet de distribution afin calculer toutes les mesures de tendance centrale (mode / médiane / moyenne) et celle de dispersion (écart-type, à calculer à $10^{-3}$ près) et de répondre à la question : quel pourcentage de l’effectif se trouve à moins d’un écart-type de la moyenne :

$$
\begin{array}{ccccc}
  32   &    31  &   29    &      28 &  34  \\
  29   &   33   &   32    &   31    &  29  \\
  29  &    29  &    34   &    33   &  31  \\ 
  31   &   31   &   31    &   33    &   30 \\ 
   29  &   32   &   33    &   34    &   30 \\
\end{array}
$$

\vspace{5mm}

\begin{tikzpicture}
	% Dimensions du repere
	\def\xmin{0} \def\xmax{15} \def\ymin{0} \def\ymax{15}
	% Grilles
%	\draw [step=0.1cm,gray,ultra thin]  (\xmin,\ymin) grid (\xmax,\ymax);
	\draw [step=0.5cm,gray, ultra thin] (\xmin,\ymin) grid (\xmax,\ymax);
	\draw [step=1cm,gray, thin] (\xmin,\ymin) grid (\xmax,\ymax);
	\draw [step=5cm,gray,thick] (\xmin,\ymin) grid (\xmax,\ymax);
\end{tikzpicture} 

\newpage

\question 

Pour son projet PPI, Justine achète à un fabriquant des coques de téléphone Univers A5 et Mono B4. Les coûts unitaires s'élèvent à CHF $1.-$ pour une coque Univers A5 et CHF $2.-$ pour un coque Mono B5 ; ces objets sont respectivement vendus CHF $5.-$ et CHF $6.-$ aux étudiants. Pour avoir assez de stock, Justine doit commander au moins $40$ coques Univers A5 et $20$ coques Mono B4. Elle peut investir au maximum que CHF $200.-$.

Sachant qu'elle peut commander au maximum $160$ coques en tout, déterminer le nombre de coques Univers A5 et Mono B4 à commander à l'optimum pour maximiser le revenu, ainsi que ce revenu maximal. 
Résoudre ce problème en dessinant le polygone de contraintes. 
Unités : un carré = 10 objets

\vspace{5mm}

\begin{tikzpicture}
	% Dimensions du repere
	\def\xmin{0} \def\xmax{15} \def\ymin{0} \def\ymax{15}
	% Grilles
%	\draw [step=0.1cm,gray,ultra thin]  (\xmin,\ymin) grid (\xmax,\ymax);
	\draw [step=0.5cm,gray, ultra thin] (\xmin,\ymin) grid (\xmax,\ymax);
	\draw [step=1cm,gray, thin] (\xmin,\ymin) grid (\xmax,\ymax);
	\draw [step=5cm,gray,thick] (\xmin,\ymin) grid (\xmax,\ymax);
\end{tikzpicture} 

\newpage

\begin{tikzpicture}
	% Dimensions du repere
	\def\xmin{0} \def\xmax{15} \def\ymin{0} \def\ymax{20}
	% Grilles
%	\draw [step=0.1cm,gray,ultra thin]  (\xmin,\ymin) grid (\xmax,\ymax);
	\draw [step=0.5cm,gray, ultra thin] (\xmin,\ymin) grid (\xmax,\ymax);
	\draw [step=1cm,gray, thin] (\xmin,\ymin) grid (\xmax,\ymax);
	\draw [step=5cm,gray,thick] (\xmin,\ymin) grid (\xmax,\ymax);
\end{tikzpicture} 

\newpage

\question Suite à une marrée noire, la surface en $m^2$ polluée est définie par la fonction $f(t)= 3000\cdot 2.7^{-0.25\cdot t}$ par rapport à $t$ le nombre de jours de "nettoyage" écoulés depuis son apparition.
(les calculs sont à arrondir à la première décimale non nulle, les réponses sont à donner à l’entier)

\begin{parts}
\part Quelle surface était polluée au moment de l'incident ?
\part Quelle surface est encore polluée le dixième jours ?
\part Quelle surface a été nettoyée uniquement le cinquième jours?
\end{parts}

\vspace{5mm}

\begin{tikzpicture}
	% Dimensions du repere
	\def\xmin{0} \def\xmax{15} \def\ymin{0} \def\ymax{15}
	% Grilles
%	\draw [step=0.1cm,gray,ultra thin]  (\xmin,\ymin) grid (\xmax,\ymax);
	\draw [step=0.5cm,gray, ultra thin] (\xmin,\ymin) grid (\xmax,\ymax);
	\draw [step=1cm,gray, thin] (\xmin,\ymin) grid (\xmax,\ymax);
	\draw [step=5cm,gray,thick] (\xmin,\ymin) grid (\xmax,\ymax);
\end{tikzpicture} 

\newpage

\question Math financières

\begin{parts}
\part Quelle est la valeur finale d’une rente de placement constituée de versements de CHF $150.-$ effectués sur un compte au taux annuel de $3\%$ à la fin de chaque mois, capitalisation mensuelle des intérêts, sur une durée de 6 ans ?

\vspace{1cm}

\begin{tikzpicture}
	% Dimensions du repere
	\def\xmin{0} \def\xmax{15} \def\ymin{0} \def\ymax{7}
	% Grilles
%	\draw [step=0.1cm,gray,ultra thin]  (\xmin,\ymin) grid (\xmax,\ymax);
	\draw [step=0.5cm,gray, ultra thin] (\xmin,\ymin) grid (\xmax,\ymax);
	\draw [step=1cm,gray, thin] (\xmin,\ymin) grid (\xmax,\ymax);
	\draw [step=5cm,gray,thick] (\xmin,\ymin) grid (\xmax,\ymax);
\end{tikzpicture} 

\part Quelle est l’annuité versée chaque fin d’année correspondant au remboursement d’un emprunt de CHF 240'000.- sur un compte au taux annuel de $12\%$, capitalisation trimestrielle des intérêts, sur une durée de 25 ans ?

\vspace{1cm}

\begin{tikzpicture}
	% Dimensions du repere
	\def\xmin{0} \def\xmax{15} \def\ymin{0} \def\ymax{7}
	% Grilles
%	\draw [step=0.1cm,gray,ultra thin]  (\xmin,\ymin) grid (\xmax,\ymax);
	\draw [step=0.5cm,gray, ultra thin] (\xmin,\ymin) grid (\xmax,\ymax);
	\draw [step=1cm,gray, thin] (\xmin,\ymin) grid (\xmax,\ymax);
	\draw [step=5cm,gray,thick] (\xmin,\ymin) grid (\xmax,\ymax);
\end{tikzpicture} 

\newpage

\part Quelle est la durée de remboursement d’une dette de CHF 1999.- s’il est effectué par des versements de CHF 300.- au début de chaque semestre sur un compte au taux annuel de $2\%$, capitalisation semestrielle, premier remboursement six mois après l'emprunt ?
\end{parts}

\vspace{1cm}

\begin{tikzpicture}
	% Dimensions du repere
	\def\xmin{0} \def\xmax{15} \def\ymin{0} \def\ymax{7}
	% Grilles
%	\draw [step=0.1cm,gray,ultra thin]  (\xmin,\ymin) grid (\xmax,\ymax);
	\draw [step=0.5cm,gray, ultra thin] (\xmin,\ymin) grid (\xmax,\ymax);
	\draw [step=1cm,gray, thin] (\xmin,\ymin) grid (\xmax,\ymax);
	\draw [step=5cm,gray,thick] (\xmin,\ymin) grid (\xmax,\ymax);
\end{tikzpicture} 

\newpage

\textbf{Précisez le numéro de l’exercice si vous utilisez cet espace réponse supplémentaire.}

\vspace{5mm}

\begin{tikzpicture}
	% Dimensions du repere
	\def\xmin{0} \def\xmax{15} \def\ymin{0} \def\ymax{19}
	% Grilles
%	\draw [step=0.1cm,gray,ultra thin]  (\xmin,\ymin) grid (\xmax,\ymax);
	\draw [step=0.5cm,gray, ultra thin] (\xmin,\ymin) grid (\xmax,\ymax);
	\draw [step=1cm,gray, thin] (\xmin,\ymin) grid (\xmax,\ymax);
	\draw [step=5cm,gray,thick] (\xmin,\ymin) grid (\xmax,\ymax);
\end{tikzpicture} 

\newpage

\begin{tikzpicture}
	% Dimensions du repere
	\def\xmin{0} \def\xmax{15} \def\ymin{0} \def\ymax{20}
	% Grilles
%	\draw [step=0.1cm,gray,ultra thin]  (\xmin,\ymin) grid (\xmax,\ymax);
	\draw [step=0.5cm,gray, ultra thin] (\xmin,\ymin) grid (\xmax,\ymax);
	\draw [step=1cm,gray, thin] (\xmin,\ymin) grid (\xmax,\ymax);
	\draw [step=5cm,gray,thick] (\xmin,\ymin) grid (\xmax,\ymax);
\end{tikzpicture} 

\newpage

\textbf{Formulaire}


\newpage



\begin{solution}
\begin{enumerate}[1.]
    \item 
            \begin{enumerate}[(a)]
                \item On peut effectuer les différentes parties de calculs :
                $$
                \left(\frac{3}{4}-\frac{7}{6} \right) = \frac{9}{12}-\frac{14}{12} = -\frac{5}{12}
                $$
                $$
                \frac{3}{5}\div 12 = \frac{3}{5}\cdot \frac{1}{12} = \frac{1}{20}
                $$
                On peut alors effectuer
                $$
                -\frac{5}{12}\cdot \frac{9}{2} = -\frac{15}{8}
                $$
                Et en mettant tout cela ensemble :
                $$
                \frac{1}{20}+\frac{15}{8} = \frac{2}{40}+\frac{75}{40} = \frac{77}{40}
                $$
                \item On commence par factoriser les différents nombres :
                $$
                108 = 2^2 \cdot 3^3, \; 8=2^3, \; 27=3^3,\; 25=5^2 \mbox{ et } 24 = 2^3\cdot 3^1
                $$
                On a donc
                \begin{align*}
                    6\sqrt{3} -65\sqrt{2}\cdot \left(3\sqrt{3}-5\right)+2\sqrt{6}&=\\
                    6\sqrt{3}-18\sqrt{6}+30\sqrt{2}+2\sqrt{6}&=\\
                    6\sqrt{3} -16\sqrt{6}+30\sqrt{2}&\\
                \end{align*}
                \item En utilisant les différentes propriétés des puissances, on a 
                $$
                x^7\cdot 2^2 \cdot x^2 \cdot 2^{-3} \cdot 3^{-3} \cdot x^{-2} \cdot -2^5 \cdot x^5 \cdot 3^{-2} \cdot x^{-2} \cdot 3^{-6} \cdot x^{-3} \cdot 2^4 \cdot 3^4 \cdot x^4
                $$
                Et donc
                $$
                =-2^8\cdot 3^{-7} \cdot x^{11} = -\frac{2^8 x^{11}}{3^7}
                $$
            \end{enumerate}
    \item 
            \begin{enumerate}[(a)]
                \item Mettons sous forme canoniques les deux équations :
                \begin{align*}
                    y-\frac{6x-3}{6} &=\frac{1}{2}+\frac{x}{6}\\
                    6y-6x+3 &=3+x\\
                    7x-6y&=0
                \end{align*}
                et
                \begin{align*}
                    30x-45y+105 &= 5x+10y-3y+3\\
                    25x-42y&=-92
                \end{align*}
                en utilisant par exemple la combinaison suivante :
                $$
                \begin{array}{llll}
                     &-25x &+42y &=92  \\
                    + & 49x & -42y &=0\\
                    \hline
                    &24x & &=92
                \end{array}
                \Rightarrow x=\frac{92}{24} = \frac{23}{6}
                $$
                En substituant :
                $$
                \frac{23}{6}\cdot 7 -6y=0 \Rightarrow y = \frac{7\cdot 23}{36}=\frac{161}{36}
                $$
                \item en faisant attention aux signes :
                \begin{align*}
                    3x-5x-10 &\geq 2x+4-\frac{x+3}{2}\\
                    6x-10x-20 &\geq 4x+8-x-3\\
                    -4x-20 &\geq 3x+5\\
                    -25 &\geq 7x\\
                    -\frac{25}{7} & \geq x
                \end{align*}
                Et donc 
                $$
                S=\left] -\infty ; -\frac{25}{7}\right[
                $$
            \end{enumerate}
    \item
            \begin{enumerate}[(a)]
                \item Notons que 
                $$
                a=-2, \; b=3, \; c=2 \mbox{ et } \Delta = 9+16=25
                $$
                On a donc d'après les formules :
                \begin{itemize}
                    \item Sommet : $\left(\frac{3}{4};\frac{25}{8}\right)$
                    \item Intersection avec l'ordonnée : $x=0 \Rightarrow y=2$
                    \item Intersections avec l'abscisse : $y=0 \Rightarrow x_1 = 2 \mbox{ et } x_2 = -\frac{1}{2}$
                \end{itemize}
                \item Par comparaison :
                $$
                -2x^2+3x+2 = 3x \ssi -2x^2 +2=0 \ssi 
                \left\{
                \begin{array}{cc}
                x_1=-1     & y_1 = -3 \\
                x_2=1     & y_2=3
                \end{array}
                \right.
                $$
                \item 
                
                \begin{center}
                \begin{tikzpicture}
	            \begin{axis}[
            	x=10mm,y=10mm,
             	axis lines=middle,
            	ymajorgrids=true,
            	xmajorgrids=true,
            	xmin=-5,
            	xmax=5,
            	ymin=-5,
            	ymax=5,
            	xtick={-5,-4,...,5},
            	ytick={-5,-4,...,5},]
            	\draw[line width=1pt,smooth,samples=100,domain=-5:5] plot(\x,{-2*(\x)^2+3*\x+2});
            	\draw[line width=1pt,smooth,samples=100,domain=-5:5] plot(\x,{3*\x});
            	\begin{scriptsize}
            	\draw[color=black] (-1,-3) node {$\times$};
            	\draw[color=black] (1,3) node {$\times$};
            	\end{scriptsize}
            	\end{axis}
            	\end{tikzpicture}
            	\end{center}
            \end{enumerate}
    \item 
            Le tableau 
            
            \begin{tabular}{llllll}
\hline
\multicolumn{1}{|l|}{caractère} & \multicolumn{1}{l|}{effectif} & \multicolumn{1}{l|}{fréquence} & \multicolumn{1}{l|}{f .cumulée} & \multicolumn{1}{l|}{f. pondérée} & \multicolumn{1}{l|}{f. p. carrée} \\ \hline
\multicolumn{1}{|l|}{28}        & \multicolumn{1}{l|}{1}        & \multicolumn{1}{l|}{0.04}      & \multicolumn{1}{l|}{0.04}       & \multicolumn{1}{l|}{1.12}        & \multicolumn{1}{l|}{31.36}        \\ \hline
\multicolumn{1}{|l|}{29}        & \multicolumn{1}{l|}{6}        & \multicolumn{1}{l|}{0.24}      & \multicolumn{1}{l|}{0.28}       & \multicolumn{1}{l|}{6.96}        & \multicolumn{1}{l|}{201.84}       \\ \hline
\multicolumn{1}{|l|}{30}        & \multicolumn{1}{l|}{2}        & \multicolumn{1}{l|}{0.08}      & \multicolumn{1}{l|}{0.36}       & \multicolumn{1}{l|}{2.4}         & \multicolumn{1}{l|}{72}           \\ \hline
\multicolumn{1}{|l|}{31}        & \multicolumn{1}{l|}{6}        & \multicolumn{1}{l|}{0.24}      & \multicolumn{1}{l|}{0.6}        & \multicolumn{1}{l|}{7.44}        & \multicolumn{1}{l|}{230.64}       \\ \hline
\multicolumn{1}{|l|}{32}        & \multicolumn{1}{l|}{3}        & \multicolumn{1}{l|}{0.12}      & \multicolumn{1}{l|}{0.72}       & \multicolumn{1}{l|}{3.84}        & \multicolumn{1}{l|}{122.88}       \\ \hline
\multicolumn{1}{|l|}{33}        & \multicolumn{1}{l|}{4}        & \multicolumn{1}{l|}{0.16}      & \multicolumn{1}{l|}{0.88}       & \multicolumn{1}{l|}{5.28}        & \multicolumn{1}{l|}{174.24}       \\ \hline
\multicolumn{1}{|l|}{34}        & \multicolumn{1}{l|}{3}        & \multicolumn{1}{l|}{0.12}      & \multicolumn{1}{l|}{1}          & \multicolumn{1}{l|}{4.08}        & \multicolumn{1}{l|}{138.72}       \\ \hline
\multicolumn{1}{|l|}{}          & \multicolumn{1}{l|}{25}       & \multicolumn{1}{l|}{1}         & \multicolumn{1}{l|}{}           & \multicolumn{1}{l|}{31.12}       & \multicolumn{1}{l|}{971.68}       \\ \hline
                                &                               &                                &                                 &                                  &                                   \\
écart type                      & 1.796                    &                                &                                 &                                  &                                  
\end{tabular}

On a ainsi : mode $=29$ et $31$, médiane $=31$, moyenne $=31.12$.

Le pourcentage recherché est dans l'intervalle $[31.12-1.796;31.12+1.796] = [29.324;32,916]$, c'est-à-dire entre $30$ et $32$ compris : $\frac{11}{25}=44\%$

\newpage

    \item 
            Commençons par le systèmes d'inéquations :
            $$
            \left\{
            \begin{array}{rl}
                1x+2y & \leq 200  \\
                1x+1y & \leq 160\\
                x & \geq 40\\
                y & \geq 20
            \end{array}
            \right.
            \mbox{ avec }
            5x+6y=\mbox{ maximum}.
            $$
            En étudiant les différents droites limites :
            \begin{align*}
            d_1 : x+2y=200 & \\
            &x=0 \; y= 100\\
            &x=200 \; y=0\\
            &\mbox{pente =}-\frac{1}{2}\\
            &\\
            d_2 : x+y=160 & \\
            &x=0 \; y= 160\\
            &x=160 \;y=0\\
            &\mbox{pente =}-1\\
            &\\
            d_{\mbox{min}} : 5x+6y=\mbox{ min} & \\
            &\mbox{pente =}-\frac{5}{6}
            \end{align*}
            On peut représenter la situation :
            \begin{center}
                \begin{tikzpicture}
	            \begin{axis}[
            	x=0.5mm,y=0.5mm,
             	axis lines=middle,
            	ymajorgrids=true,
            	xmajorgrids=true,
            	xmin=0,
            	xmax=250,
            	ymin=0,
            	ymax=200,
            	xtick={0,25,...,250},
            	ytick={0,25,...,200},]
            	\draw[line width=1pt,smooth,samples=100,domain=0:250] plot(\x,{-\x/2+100});
            	\draw[line width=1pt,smooth,samples=100,domain=0:250] plot(\x,{-\x+160});
            	\draw[line width=1pt,smooth,samples=100,domain=0:250] plot(\x,{20});
            	\draw[line width=1pt,smooth,samples=100,domain=0:250] plot(\x,{-\x+160});
            	\draw[line width=1pt,smooth,samples=100,domain=0:250, dashed] plot(\x,{(-5*\x+840)/6});
            	\draw[line width=1pt,smooth,samples=100,domain=0:200] plot[variable=\t] (40,\t);
            	\fill[pattern=north east lines]
                (40,20) -- (140,20) 
                -- (120,40)
                -- (40,80) -- cycle;
                \draw[color=black] (50,100) node[below] {$d_m$};
                \draw[color=black] (120,40) node[above right] {Maximum = $(120, 40)$};
                \draw (120,40)  node {$\bullet$};
            	\end{axis}
            	\end{tikzpicture}
            	\end{center}
            	On détermine que le point optimum se trouve aux intersectcions de $d_1$ et $d_2$ :
            	$$
            	\left\{
            	\begin{array}{ll}
            	   x+2y  &=200  \\
            	    x+y &=100 
            	\end{array}
            	\right.
            	\Rightarrow
            	\left\{
            	\begin{array}{ll}
            	   x  &=120  \\
            	   y &=40 
            	\end{array}
            	\right.
            	$$
            	Le revenu est donc de $5\cdot 120 + 6\cdot 40 = 840$.
    \item 
            \begin{enumerate}[(a)]
                \item 
                $$
                f(0) = 3000\cdot 2.7^{-0.25\cdot 0} = 3000
                $$
                \item 
                $$
                f(10) = 3000\cdot 2.7^{-0.25\cdot 10} = 250
                $$
                \item 
                $$
                f(4)-f(5) = 3000\cdot 2.7^{-0.25\cdot 4}-3000\cdot 2.7^{-0.25\cdot 5} = 244
                $$
            \end{enumerate}
    \item 
            \begin{enumerate}[(a)]
            \item 
            $$
            150 \cdot \frac{1.0025^{72}-1}{0.0025} = 11'816.91
            $$
            \item
            $$
            240'000=a\cdot \frac{1-1.03^{-100}}{0.03} \Rightarrow a=759.52
            $$
            \item 
            $$
            1999 = 300\cdot \frac{1-1.01^{-n}}{0.01} \Rightarrow n\simeq 6.93
            $$
            Donc $6$ remboursements entiers et $1$ remboursement partiel.
            \end{enumerate}
\end{enumerate}


\end{solution}
\end{questions}
\end{document}


\end{document}